\section{Commutator and return receipts}

Define the encounter commutator
\[
\comm=[A,B]=ABA^{-1}B^{-1}.
\]
Because (A) and (B) are involutions,
\[
\comm=ABAB=H^2.
\]
Since (H) has order (2n), the element (H^2) has order (n). It records
the difference between the forward product (AB) and reverse product (BA):
\[
AB=\comm\,BA.
\]

The commutator subgroup of the dihedral completion is
\[
[\langle A,B\rangle,\langle A,B\rangle]
=\langle H^2\rangle
=\langle\comm\rangle.
\]
Indeed, quotienting by (langle H^2\rangle) makes the images of (A) and
(B) commuting involutions, while the commutator calculation shows that
(H^2) is itself generated by a commutator.

Now define the binary return receipt
\[
\receipt=H^n.
\]
It is the unique nonidentity element of order two in the cyclic subgroup
(langle H\rangle\cong C_{2n}). Under the orientation projection
\[
\langle H\rangle\longrightarrow C_n,
\qquad H\longmapsto r,
\]
the kernel is exactly
\[
\{1,\receipt\}.
\]

The central question is whether the deck receipt (
eceipt) belongs to the
commutator tower (langle\comm\rangle). This asks whether the congruence
\[
2k\equiv n\pmod{2n}
\]
has a solution.

